% Seminar 2 - Pagina 4. Fragment pentru % Seminar 2 - Pagina 4. Fragment pentru % Seminar 2 - Pagina 4. Fragment pentru % Seminar 2 - Pagina 4. Fragment pentru \input{seminars/pages/seminar2-pagina4-ro}.
% Fara preambul si fara comanda de schimbare fortata a paginii.
% NOTE EDITORIALE:
% - Pagina continua demonstratia inceputa la punctul c) pe pagina 3.
%   Manuscrisul revine mai intai la demonstratiile a) si b), apoi la c).
% - La punctul b), egalitatea in cazul f injectiva este demonstrata
%   cu antecedente a_i care au aceeasi imagine; se presupune o familie
%   indexata nevida, fara a explicita aceasta conditie.
% - Exemplul cu A={-1,0,1} si f(x)=x^2 pare sa scrie f(X_1)=A,
%   desi matematic f({-1,1})={1}. Am pastrat citirea provizorie A.
% - Notele din dreapta despre imaginile unor intervale par sa foloseasca
%   tot f(x)=x^2, dar domeniul nu este precizat acolo. In manuscris
%   f([-1,3]) pare notata [1,9], care ar trebui sa fie [0,9].
%   Am pastrat expresia citita si am semnalat aceasta neconcordanta.
% - Calculul din chenar are semne de interval greu de distins;
%   semnele [ ] de mai jos sunt lecturi provizorii.
% - Ultimul rand al punctului c) se intrerupe in manuscris;
%   nu este completat argumentul si nici punctul d).

\subsection*{Seminar 2 - Pagina 4}

\noindent\underline{Sol. (continuare)}

\medskip
\noindent $a)$ \underline{$\subseteq$}.
Fie $y\in f\left(\bigcup_{i\in I}X_i\right)$. Atunci
\[
\exists\,x\in\bigcup_{i\in I}X_i\quad\text{a.î.}\quad f(x)=y.
\]
\noindent Rezultă că există $i_0\in I$ cu $x\in X_{i_0}$ şi $f(x)=y$, deci
\[
 y\in f(X_{i_0})\subseteq\bigcup_{i\in I}f(X_i).
\]
\noindent \underline{$\supseteq$}.
Există $i_0\in I$ a.î.
$y\in f(X_{i_0})$, de unde
\[
 y\in f\left(\bigcup_{i\in I}X_i\right).
\]

\medskip
\noindent $b)$ Fie
\[
 y\in f\left(\bigcap_{i\in I}X_i\right).
\]
\noindent Există $x\in\bigcap_{i\in I}X_i$ cu $f(x)=y$, deci
\[
 y\in\bigcap_{i\in I}f(X_i).
\]
\noindent P.p.\ $f$ injectivă. Dacă
$y\in\bigcap_{i\in I}f(X_i)$, atunci
$y\in f(X_i)$ pentru orice $i\in I$. Prin urmare,
\[
 (\forall i\in I)(\exists a_i\in X_i):\ f(a_i)=y.
\]
\noindent Din injectivitate, $a_i=a_j$ pentru orice $i,j\in I$.
Notăm acest element comun cu $a$. Atunci
\[
 f(a)=y,\qquad a\in\bigcap_{i\in I}X_i,
 \qquad y\in f\left(\bigcap_{i\in I}X_i\right).
\]

\medskip
\noindent\underline{Ex.} În manuscris apare şi exemplul
\[
 A=\{-1,0,1\},\qquad B=\{0,1\},\qquad f(x)=x^2,
\]
\[
 X_1=\{-1,1\},\qquad X_2=\{-1,0\},
\]
\[
 f(X_1)=A,\qquad f(X_2)=\{0,1\}.
\]

% Schita marginala a parabolei y=x^2.
\begin{center}
\begin{tikzpicture}[>=stealth,scale=.77]
  \draw[->] (-1.8,0)--(3.0,0) node[right] {$x$};
  \draw[->] (0,-.25)--(0,3.45) node[above] {$y$};
  \foreach \x in {-1,0,1,2,3}{\draw (\x,.06)--(\x,-.06) node[below] {\scriptsize $\x$};}
  \draw[domain=-1.35:2.35,samples=70,smooth,thick] plot (\x,{.5*\x*\x});
\end{tikzpicture}
\end{center}

% Exemple marginale: semnele de interval si unele limite sunt greu lizibile.
\noindent\underline{Calcule din margine (lecturi provizorii):}
\[
\begin{aligned}
 f([1,2])&=[1,4], & f([-1,3])&=[1,9],\\
 f([1,2])&=f([1,3]\cap[1,2])=[1,4].
\end{aligned}
\]
\[
\boxed{\begin{aligned}
 f([-1,0])&=[0,1],\\
 f([0,2])&=[0,4],\\
 f([-1,0]\cap[0,2])&=f(0)=0,\\
 f([-1,0])\cap f([0,2])&=[0,1].
\end{aligned}}
\]

\medskip
\noindent $c)$ \underline{$\subseteq$}.
Fie
\[
 x\in f^{-1}\left(\bigcup_{j\in J}Y_j\right).
\]
\noindent Atunci
\[
 f(x)\in\bigcup_{j\in J}Y_j
 \quad\Longrightarrow\quad
 (\exists j_0\in J):\ f(x)\in Y_{j_0}
 \quad\Longrightarrow\quad
 x\in f^{-1}(Y_{j_0})
 \subseteq\bigcup_{j\in J}f^{-1}(Y_j).
\]
\noindent \underline{$\supseteq$}.
Dacă $x\in\bigcup_{j\in J}f^{-1}(Y_j)$, există $j_0\in J$
astfel încât $x\in f^{-1}(Y_{j_0})$, de unde
\[
 f(x)\in Y_{j_0}\ \Longrightarrow\ f(x)\in\bigcup_{j\in J}Y_j.
\]
% Sfarsitul paginii: argumentul nu este dus mai departe.
.
% Fara preambul si fara comanda de schimbare fortata a paginii.
% NOTE EDITORIALE:
% - Pagina continua demonstratia inceputa la punctul c) pe pagina 3.
%   Manuscrisul revine mai intai la demonstratiile a) si b), apoi la c).
% - La punctul b), egalitatea in cazul f injectiva este demonstrata
%   cu antecedente a_i care au aceeasi imagine; se presupune o familie
%   indexata nevida, fara a explicita aceasta conditie.
% - Exemplul cu A={-1,0,1} si f(x)=x^2 pare sa scrie f(X_1)=A,
%   desi matematic f({-1,1})={1}. Am pastrat citirea provizorie A.
% - Notele din dreapta despre imaginile unor intervale par sa foloseasca
%   tot f(x)=x^2, dar domeniul nu este precizat acolo. In manuscris
%   f([-1,3]) pare notata [1,9], care ar trebui sa fie [0,9].
%   Am pastrat expresia citita si am semnalat aceasta neconcordanta.
% - Calculul din chenar are semne de interval greu de distins;
%   semnele [ ] de mai jos sunt lecturi provizorii.
% - Ultimul rand al punctului c) se intrerupe in manuscris;
%   nu este completat argumentul si nici punctul d).

\subsection*{Seminar 2 - Pagina 4}

\noindent\underline{Sol. (continuare)}

\medskip
\noindent $a)$ \underline{$\subseteq$}.
Fie $y\in f\left(\bigcup_{i\in I}X_i\right)$. Atunci
\[
\exists\,x\in\bigcup_{i\in I}X_i\quad\text{a.î.}\quad f(x)=y.
\]
\noindent Rezultă că există $i_0\in I$ cu $x\in X_{i_0}$ şi $f(x)=y$, deci
\[
 y\in f(X_{i_0})\subseteq\bigcup_{i\in I}f(X_i).
\]
\noindent \underline{$\supseteq$}.
Există $i_0\in I$ a.î.
$y\in f(X_{i_0})$, de unde
\[
 y\in f\left(\bigcup_{i\in I}X_i\right).
\]

\medskip
\noindent $b)$ Fie
\[
 y\in f\left(\bigcap_{i\in I}X_i\right).
\]
\noindent Există $x\in\bigcap_{i\in I}X_i$ cu $f(x)=y$, deci
\[
 y\in\bigcap_{i\in I}f(X_i).
\]
\noindent P.p.\ $f$ injectivă. Dacă
$y\in\bigcap_{i\in I}f(X_i)$, atunci
$y\in f(X_i)$ pentru orice $i\in I$. Prin urmare,
\[
 (\forall i\in I)(\exists a_i\in X_i):\ f(a_i)=y.
\]
\noindent Din injectivitate, $a_i=a_j$ pentru orice $i,j\in I$.
Notăm acest element comun cu $a$. Atunci
\[
 f(a)=y,\qquad a\in\bigcap_{i\in I}X_i,
 \qquad y\in f\left(\bigcap_{i\in I}X_i\right).
\]

\medskip
\noindent\underline{Ex.} În manuscris apare şi exemplul
\[
 A=\{-1,0,1\},\qquad B=\{0,1\},\qquad f(x)=x^2,
\]
\[
 X_1=\{-1,1\},\qquad X_2=\{-1,0\},
\]
\[
 f(X_1)=A,\qquad f(X_2)=\{0,1\}.
\]

% Schita marginala a parabolei y=x^2.
\begin{center}
\begin{tikzpicture}[>=stealth,scale=.77]
  \draw[->] (-1.8,0)--(3.0,0) node[right] {$x$};
  \draw[->] (0,-.25)--(0,3.45) node[above] {$y$};
  \foreach \x in {-1,0,1,2,3}{\draw (\x,.06)--(\x,-.06) node[below] {\scriptsize $\x$};}
  \draw[domain=-1.35:2.35,samples=70,smooth,thick] plot (\x,{.5*\x*\x});
\end{tikzpicture}
\end{center}

% Exemple marginale: semnele de interval si unele limite sunt greu lizibile.
\noindent\underline{Calcule din margine (lecturi provizorii):}
\[
\begin{aligned}
 f([1,2])&=[1,4], & f([-1,3])&=[1,9],\\
 f([1,2])&=f([1,3]\cap[1,2])=[1,4].
\end{aligned}
\]
\[
\boxed{\begin{aligned}
 f([-1,0])&=[0,1],\\
 f([0,2])&=[0,4],\\
 f([-1,0]\cap[0,2])&=f(0)=0,\\
 f([-1,0])\cap f([0,2])&=[0,1].
\end{aligned}}
\]

\medskip
\noindent $c)$ \underline{$\subseteq$}.
Fie
\[
 x\in f^{-1}\left(\bigcup_{j\in J}Y_j\right).
\]
\noindent Atunci
\[
 f(x)\in\bigcup_{j\in J}Y_j
 \quad\Longrightarrow\quad
 (\exists j_0\in J):\ f(x)\in Y_{j_0}
 \quad\Longrightarrow\quad
 x\in f^{-1}(Y_{j_0})
 \subseteq\bigcup_{j\in J}f^{-1}(Y_j).
\]
\noindent \underline{$\supseteq$}.
Dacă $x\in\bigcup_{j\in J}f^{-1}(Y_j)$, există $j_0\in J$
astfel încât $x\in f^{-1}(Y_{j_0})$, de unde
\[
 f(x)\in Y_{j_0}\ \Longrightarrow\ f(x)\in\bigcup_{j\in J}Y_j.
\]
% Sfarsitul paginii: argumentul nu este dus mai departe.
.
% Fara preambul si fara comanda de schimbare fortata a paginii.
% NOTE EDITORIALE:
% - Pagina continua demonstratia inceputa la punctul c) pe pagina 3.
%   Manuscrisul revine mai intai la demonstratiile a) si b), apoi la c).
% - La punctul b), egalitatea in cazul f injectiva este demonstrata
%   cu antecedente a_i care au aceeasi imagine; se presupune o familie
%   indexata nevida, fara a explicita aceasta conditie.
% - Exemplul cu A={-1,0,1} si f(x)=x^2 pare sa scrie f(X_1)=A,
%   desi matematic f({-1,1})={1}. Am pastrat citirea provizorie A.
% - Notele din dreapta despre imaginile unor intervale par sa foloseasca
%   tot f(x)=x^2, dar domeniul nu este precizat acolo. In manuscris
%   f([-1,3]) pare notata [1,9], care ar trebui sa fie [0,9].
%   Am pastrat expresia citita si am semnalat aceasta neconcordanta.
% - Calculul din chenar are semne de interval greu de distins;
%   semnele [ ] de mai jos sunt lecturi provizorii.
% - Ultimul rand al punctului c) se intrerupe in manuscris;
%   nu este completat argumentul si nici punctul d).

\subsection*{Seminar 2 - Pagina 4}

\noindent\underline{Sol. (continuare)}

\medskip
\noindent $a)$ \underline{$\subseteq$}.
Fie $y\in f\left(\bigcup_{i\in I}X_i\right)$. Atunci
\[
\exists\,x\in\bigcup_{i\in I}X_i\quad\text{a.î.}\quad f(x)=y.
\]
\noindent Rezultă că există $i_0\in I$ cu $x\in X_{i_0}$ şi $f(x)=y$, deci
\[
 y\in f(X_{i_0})\subseteq\bigcup_{i\in I}f(X_i).
\]
\noindent \underline{$\supseteq$}.
Există $i_0\in I$ a.î.
$y\in f(X_{i_0})$, de unde
\[
 y\in f\left(\bigcup_{i\in I}X_i\right).
\]

\medskip
\noindent $b)$ Fie
\[
 y\in f\left(\bigcap_{i\in I}X_i\right).
\]
\noindent Există $x\in\bigcap_{i\in I}X_i$ cu $f(x)=y$, deci
\[
 y\in\bigcap_{i\in I}f(X_i).
\]
\noindent P.p.\ $f$ injectivă. Dacă
$y\in\bigcap_{i\in I}f(X_i)$, atunci
$y\in f(X_i)$ pentru orice $i\in I$. Prin urmare,
\[
 (\forall i\in I)(\exists a_i\in X_i):\ f(a_i)=y.
\]
\noindent Din injectivitate, $a_i=a_j$ pentru orice $i,j\in I$.
Notăm acest element comun cu $a$. Atunci
\[
 f(a)=y,\qquad a\in\bigcap_{i\in I}X_i,
 \qquad y\in f\left(\bigcap_{i\in I}X_i\right).
\]

\medskip
\noindent\underline{Ex.} În manuscris apare şi exemplul
\[
 A=\{-1,0,1\},\qquad B=\{0,1\},\qquad f(x)=x^2,
\]
\[
 X_1=\{-1,1\},\qquad X_2=\{-1,0\},
\]
\[
 f(X_1)=A,\qquad f(X_2)=\{0,1\}.
\]

% Schita marginala a parabolei y=x^2.
\begin{center}
\begin{tikzpicture}[>=stealth,scale=.77]
  \draw[->] (-1.8,0)--(3.0,0) node[right] {$x$};
  \draw[->] (0,-.25)--(0,3.45) node[above] {$y$};
  \foreach \x in {-1,0,1,2,3}{\draw (\x,.06)--(\x,-.06) node[below] {\scriptsize $\x$};}
  \draw[domain=-1.35:2.35,samples=70,smooth,thick] plot (\x,{.5*\x*\x});
\end{tikzpicture}
\end{center}

% Exemple marginale: semnele de interval si unele limite sunt greu lizibile.
\noindent\underline{Calcule din margine (lecturi provizorii):}
\[
\begin{aligned}
 f([1,2])&=[1,4], & f([-1,3])&=[1,9],\\
 f([1,2])&=f([1,3]\cap[1,2])=[1,4].
\end{aligned}
\]
\[
\boxed{\begin{aligned}
 f([-1,0])&=[0,1],\\
 f([0,2])&=[0,4],\\
 f([-1,0]\cap[0,2])&=f(0)=0,\\
 f([-1,0])\cap f([0,2])&=[0,1].
\end{aligned}}
\]

\medskip
\noindent $c)$ \underline{$\subseteq$}.
Fie
\[
 x\in f^{-1}\left(\bigcup_{j\in J}Y_j\right).
\]
\noindent Atunci
\[
 f(x)\in\bigcup_{j\in J}Y_j
 \quad\Longrightarrow\quad
 (\exists j_0\in J):\ f(x)\in Y_{j_0}
 \quad\Longrightarrow\quad
 x\in f^{-1}(Y_{j_0})
 \subseteq\bigcup_{j\in J}f^{-1}(Y_j).
\]
\noindent \underline{$\supseteq$}.
Dacă $x\in\bigcup_{j\in J}f^{-1}(Y_j)$, există $j_0\in J$
astfel încât $x\in f^{-1}(Y_{j_0})$, de unde
\[
 f(x)\in Y_{j_0}\ \Longrightarrow\ f(x)\in\bigcup_{j\in J}Y_j.
\]
% Sfarsitul paginii: argumentul nu este dus mai departe.
.
% Fara preambul si fara comanda de schimbare fortata a paginii.
% NOTE EDITORIALE:
% - Pagina continua demonstratia inceputa la punctul c) pe pagina 3.
%   Manuscrisul revine mai intai la demonstratiile a) si b), apoi la c).
% - La punctul b), egalitatea in cazul f injectiva este demonstrata
%   cu antecedente a_i care au aceeasi imagine; se presupune o familie
%   indexata nevida, fara a explicita aceasta conditie.
% - Exemplul cu A={-1,0,1} si f(x)=x^2 pare sa scrie f(X_1)=A,
%   desi matematic f({-1,1})={1}. Am pastrat citirea provizorie A.
% - Notele din dreapta despre imaginile unor intervale par sa foloseasca
%   tot f(x)=x^2, dar domeniul nu este precizat acolo. In manuscris
%   f([-1,3]) pare notata [1,9], care ar trebui sa fie [0,9].
%   Am pastrat expresia citita si am semnalat aceasta neconcordanta.
% - Calculul din chenar are semne de interval greu de distins;
%   semnele [ ] de mai jos sunt lecturi provizorii.
% - Ultimul rand al punctului c) se intrerupe in manuscris;
%   nu este completat argumentul si nici punctul d).

\subsection*{Seminar 2 - Pagina 4}

\noindent\underline{Sol. (continuare)}

\medskip
\noindent $a)$ \underline{$\subseteq$}.
Fie $y\in f\left(\bigcup_{i\in I}X_i\right)$. Atunci
\[
\exists\,x\in\bigcup_{i\in I}X_i\quad\text{a.î.}\quad f(x)=y.
\]
\noindent Rezultă că există $i_0\in I$ cu $x\in X_{i_0}$ şi $f(x)=y$, deci
\[
 y\in f(X_{i_0})\subseteq\bigcup_{i\in I}f(X_i).
\]
\noindent \underline{$\supseteq$}.
Există $i_0\in I$ a.î.
$y\in f(X_{i_0})$, de unde
\[
 y\in f\left(\bigcup_{i\in I}X_i\right).
\]

\medskip
\noindent $b)$ Fie
\[
 y\in f\left(\bigcap_{i\in I}X_i\right).
\]
\noindent Există $x\in\bigcap_{i\in I}X_i$ cu $f(x)=y$, deci
\[
 y\in\bigcap_{i\in I}f(X_i).
\]
\noindent P.p.\ $f$ injectivă. Dacă
$y\in\bigcap_{i\in I}f(X_i)$, atunci
$y\in f(X_i)$ pentru orice $i\in I$. Prin urmare,
\[
 (\forall i\in I)(\exists a_i\in X_i):\ f(a_i)=y.
\]
\noindent Din injectivitate, $a_i=a_j$ pentru orice $i,j\in I$.
Notăm acest element comun cu $a$. Atunci
\[
 f(a)=y,\qquad a\in\bigcap_{i\in I}X_i,
 \qquad y\in f\left(\bigcap_{i\in I}X_i\right).
\]

\medskip
\noindent\underline{Ex.} În manuscris apare şi exemplul
\[
 A=\{-1,0,1\},\qquad B=\{0,1\},\qquad f(x)=x^2,
\]
\[
 X_1=\{-1,1\},\qquad X_2=\{-1,0\},
\]
\[
 f(X_1)=A,\qquad f(X_2)=\{0,1\}.
\]

% Schita marginala a parabolei y=x^2.
\begin{center}
\begin{tikzpicture}[>=stealth,scale=.77]
  \draw[->] (-1.8,0)--(3.0,0) node[right] {$x$};
  \draw[->] (0,-.25)--(0,3.45) node[above] {$y$};
  \foreach \x in {-1,0,1,2,3}{\draw (\x,.06)--(\x,-.06) node[below] {\scriptsize $\x$};}
  \draw[domain=-1.35:2.35,samples=70,smooth,thick] plot (\x,{.5*\x*\x});
\end{tikzpicture}
\end{center}

% Exemple marginale: semnele de interval si unele limite sunt greu lizibile.
\noindent\underline{Calcule din margine (lecturi provizorii):}
\[
\begin{aligned}
 f([1,2])&=[1,4], & f([-1,3])&=[1,9],\\
 f([1,2])&=f([1,3]\cap[1,2])=[1,4].
\end{aligned}
\]
\[
\boxed{\begin{aligned}
 f([-1,0])&=[0,1],\\
 f([0,2])&=[0,4],\\
 f([-1,0]\cap[0,2])&=f(0)=0,\\
 f([-1,0])\cap f([0,2])&=[0,1].
\end{aligned}}
\]

\medskip
\noindent $c)$ \underline{$\subseteq$}.
Fie
\[
 x\in f^{-1}\left(\bigcup_{j\in J}Y_j\right).
\]
\noindent Atunci
\[
 f(x)\in\bigcup_{j\in J}Y_j
 \quad\Longrightarrow\quad
 (\exists j_0\in J):\ f(x)\in Y_{j_0}
 \quad\Longrightarrow\quad
 x\in f^{-1}(Y_{j_0})
 \subseteq\bigcup_{j\in J}f^{-1}(Y_j).
\]
\noindent \underline{$\supseteq$}.
Dacă $x\in\bigcup_{j\in J}f^{-1}(Y_j)$, există $j_0\in J$
astfel încât $x\in f^{-1}(Y_{j_0})$, de unde
\[
 f(x)\in Y_{j_0}\ \Longrightarrow\ f(x)\in\bigcup_{j\in J}Y_j.
\]
% Sfarsitul paginii: argumentul nu este dus mai departe.
